\chapter{Checklist Técnico de Fechamento}

\section{Arquivos e versão}

\begin{itemize}[label=$\square$]
  \item Confirme que o trabalho usa uma cópia íntegra do UnDF-TeX 0.1.0.
  \item Confira os arquivos de entrada, capítulos, imagens e bibliografia.
  \item Remova apenas artefatos temporários; preserve fontes e arquivos necessários à reprodução.
  \item Verifique se não há dados de exemplo remanescentes.
\end{itemize}

\section{Compilação}

\begin{itemize}[label=$\square$]
  \item Compile com LuaLaTeX e Biber pelo fluxo recomendado.
  \item Repita as passagens até sumário, listas e rótulos estabilizarem.
  \item Corrija todos os erros fatais e avalie cada aviso relevante.
  \item Confirme que não há citações ou referências cruzadas indefinidas.
\end{itemize}

\section{Conferência visual}

\begin{itemize}[label=$\square$]
  \item Abra o PDF gerado, não apenas a visualização do editor.
  \item Confira capa, folha de rosto, aprovação, resumos e sumário.
  \item Examine transições de página, objetos extensos e páginas com pouco conteúdo.
  \item Confira nitidez, legenda, fonte e chamada de cada objeto.
  \item Verifique links, marcadores, numeração e páginas incluídas externamente.
  \item Confirme ordem de leitura, contraste, legibilidade e alternativas textuais quando o fluxo de produção as suportar.
  \item Faça uma leitura em escala de 100\% e outra em visão geral de miniaturas.
\end{itemize}

\section{Entrega técnica}

\begin{itemize}[label=$\square$]
  \item Gere o PDF a partir da versão efetivamente aprovada para entrega.
  \item Use nome de arquivo claro, sem marcas como \texttt{final-final}.
  \item Confirme abertura em outro leitor de PDF.
  \item Arquive o código-fonte, a bibliografia e os recursos usados na geração.
\end{itemize}
